\begin{figure}[htbp]
\centering
\begin{tikzpicture}[schema base]
\node[schema titolo] at (0,0.4) {Pagina sorgente};
\node[schema titolo] at (6.3,0.4) {Contenuti organizzati nel manuale};
\path[schema pannello] (0,0) rectangle (4.5,-7.35);
\node[schema nota] at (2.25,-0.4) {Pagina fisica $p$};
\path[fill=schemaCartaBlu] (0.3,-0.8) rectangle (4.2,-3.9);
\node[anchor=west,font=\schemafont\small\bfseries] at (0.5,-1.15) {A. Comandi};
\foreach \riga in {-1.65,-1.9,-2.15} {\draw[schemaGrigio!70] (0.5,\riga) -- (3.9,\riga);}
\pic[scale=1] at (2.25,-3) {immagine};
\draw[schemaVerde,dashed] (1.35,-3.7) rectangle (3.15,-2.3);
\path[fill=schemaCartaGialla!60] (0.3,-4.2) rectangle (4.2,-6.95);
\node[anchor=west,font=\schemafont\small\bfseries] at (0.5,-4.6) {B. Parametri};
\pic[scale=1.3] at (2.25,-5.8) {tabella};
\node[schema nota,text width=4.4cm] at (2.25,-8.05)
  {Due sezioni iniziano sulla stessa pagina: servono titoli e ordine di lettura.};
\draw[schema freccia] (4.7,-3.6) -- (6,-3.6);
\node[schema nota,rotate=90] at (5.35,-5.4) {Assemblaggio};
\node[schema tesi,text width=5.9cm,minimum height=0.8cm] (manuale) at (11.2,-0.5)
  {\texttt{ManualDocument}};
\node[schema blocco,text width=3.6cm,minimum height=0.75cm] (capitoloa) at (10,-2)
  {\texttt{Chapter}: A. Comandi};
\node[schema esterno,text width=3.6cm,minimum height=0.75cm] (capitolob) at (10,-5.4)
  {\texttt{Chapter}: B. Parametri};
\draw[schema icona] (manuale.west) -- (6.6,-0.5) -- (6.6,-5.4) -- (capitolob.west);
\draw[schema icona] (6.6,-2) -- (capitoloa.west);
\pic[scale=0.35] at (10,-3.05) {testo};
\node[schema etichetta,anchor=west] at (10.5,-3.05) {Testo};
\pic[scale=0.4] at (10,-4.15) {immagine};
\node[schema etichetta,anchor=west] at (10.5,-4.15) {Immagine + asset};
\draw[schema icona] (capitoloa.south) -- (10,-2.55) -- (9.1,-2.55) -- (9.1,-4.15) -- (9.6,-4.15);
\draw[schema icona] (9.1,-3.05) -- (9.6,-3.05);
\pic[scale=0.4] at (10,-6.6) {tabella};
\node[schema etichetta,anchor=west] at (10.5,-6.6) {Tabella};
\draw[schema icona] (capitolob.south) -- (10,-5.98) -- (9.1,-5.98) -- (9.1,-6.6) -- (9.6,-6.6);
\node[schema nota,text width=8.3cm] at (10.5,-8.05)
  {La gerarchia può includere sottocapitoli; gli elementi precedenti alla prima sezione restano alla radice.};
\path[schema pannello,fill=schemaCartaVerde!35] (0,-9) rectangle (15,-11.3);
\pic[scale=0.6] at (0.9,-10.15) {immagine};
\node[schema etichetta,text width=3.4cm] at (3.3,-10.1)
  {\textbf{Elemento immagine}\\Ritaglio conservato\\e riferimento relativo};
\draw[schema freccia,draw=schemaVerde] (5.2,-10.1) -- (6,-10.1);
\node[schema etichetta,text width=3.6cm] at (8,-10.1)
  {\textbf{Fonte e collocazione}\\Pagina, bounding box\\e capitolo padre};
\node[schema etichetta,text width=4cm] at (12.6,-10.1)
  {\textbf{Arricchimento distinto}\\Caption originale / generata\\e provenienza del modello};
\end{tikzpicture}
\caption{Esempio illustrativo di assemblaggio e tracciabilità, non estratto da uno dei manuali valutati. Sono rappresentati solo alcuni elementi figli. I titoli delimitano sezioni diverse anche sulla stessa pagina; il ritaglio rimane collegato alla fonte e distinto dalla descrizione generata. I dati geometrici e di provenienza sono disponibili secondo il percorso di acquisizione e l'arricchimento eseguito.}
\label{fig:gerarchia-provenienza}
\end{figure}